% Entraînement — Réduction (Banque PT 2018 et 2015, Maths A) — corrigé Le Collé (non officiel)
\documentclass[a4paper,11pt]{article}

% ============================================================
%  Le Collé — Algèbre linéaire (compléments de deuxième année)
%  © 2026 Ahmed Amine Ouqlani
% ============================================================

% ---------- Polices ----------
\usepackage[utf8]{inputenc}
\usepackage[T1]{fontenc}
\usepackage[sc]{mathpazo}          % texte et mathématiques : Palatino
\usepackage[scaled=0.88]{helvet}   % titres et repères : linéale
\linespread{1.07}
\usepackage{amsmath,amssymb,amsthm,mathrsfs,stmaryrd}
\usepackage{microtype}

% ---------- Mise en page ----------
\usepackage[a4paper,top=2.3cm,bottom=2.3cm,left=2.4cm,right=2.4cm,
            headheight=14pt,headsep=0.8cm,footskip=1.1cm]{geometry}
\setlength{\parindent}{0pt}
\setlength{\parskip}{0.45em plus 0.15em minus 0.1em}
\raggedbottom
\emergencystretch=1.5em
\clubpenalty=10000 \widowpenalty=10000 \displaywidowpenalty=10000 \brokenpenalty=10000

% ---------- Couleurs : un bleu de nuit, un ocre, des gris ----------
\usepackage{xcolor}
\definecolor{Ink}{HTML}{1C2533}
\definecolor{Bleu}{HTML}{1F3A5F}
\definecolor{BleuMoyen}{HTML}{4A6A8F}
\definecolor{BleuClair}{HTML}{9DB2C9}
\definecolor{Fond}{HTML}{F2F4F7}
\definecolor{Ocre}{HTML}{B07A26}
\definecolor{Brique}{HTML}{A23B2A}
\definecolor{Gris}{HTML}{5E6672}
\definecolor{Filet}{HTML}{C5CCD6}

% ---------- Packages ----------
\usepackage{enumitem}
\setlist{itemsep=0.25em,topsep=0.35em,parsep=0pt}
\setlist[itemize]{label=\textcolor{BleuMoyen}{\small\textbullet},leftmargin=1.4em}
\setlist[enumerate]{leftmargin=1.8em,label=\textsf{\small\color{Bleu}\arabic*.}}
\usepackage{booktabs,array,multicol}
\usepackage{needspace,float}
\usepackage[font=small,labelfont={sf,color=Bleu},labelsep=period,skip=6pt]{caption}
\usepackage{tikz}
\usetikzlibrary{arrows.meta,calc,decorations.markings}
\usepackage{pgfplots}
\pgfplotsset{compat=1.17}
\usepackage{tkz-tab}
\usepackage[most]{tcolorbox}
\usepackage{titlesec,titletoc}
\usepackage{fancyhdr,lastpage,etoolbox}
\usepackage[colorlinks=true,linkcolor=Bleu,urlcolor=Bleu,citecolor=Bleu,
            pdftitle={Entraînement — Réduction (Banque PT 2018, Maths A)},pdfauthor={Ahmed Amine Ouqlani}]{hyperref}

\numberwithin{equation}{section}
\numberwithin{figure}{section}
\numberwithin{table}{section}

% ---------- Titres ----------
% numéro du chapitre dans le parcours PT (affiché), distinct de la numérotation interne
\newcommand{\grandnum}{\thesection}\newcommand{\ongletcouleur}{Bleu}
% Chapitres (sections) : grand numéro, titre en linéale, filet ocre
\titleformat{\section}[display]
  {\normalfont\sffamily\color{Ink}}
  {{\small\color{Gris}\textls[200]{\MakeUppercase{Chapitre}}}\\[-0.2em]%
   {\fontsize{64}{64}\selectfont\color{BleuClair}\grandnum}}
  {0.4em}
  {\fontsize{25}{30}\selectfont\bfseries}
  [\vspace{0.5em}{\color{Ocre}\rule{3.2cm}{1.4pt}}]
\titlespacing*{\section}{0pt}{0pt}{1.6em}
\titleformat{name=\section,numberless}[display]
  {\normalfont\sffamily\color{Ink}}{}{0pt}{\fontsize{25}{30}\selectfont\bfseries}
  [\vspace{0.5em}{\color{Ocre}\rule{3.2cm}{1.4pt}}]
\titlespacing*{name=\section,numberless}{0pt}{0pt}{1.6em}

\titleformat{\subsection}
  {\needspace{5\baselineskip}\normalfont\sffamily\Large\bfseries\color{Ink}}
  {\color{Bleu}\thesubsection}{0.8em}{}
\titlespacing*{\subsection}{0pt}{1.8em}{0.5em}
\titleformat{\subsubsection}
  {\needspace{5\baselineskip}\normalfont\sffamily\large\bfseries\color{Ink}}
  {\color{Bleu}\thesubsubsection}{0.7em}{}
\titlespacing*{\subsubsection}{0pt}{1.6em}{0.4em}
\titleformat{\paragraph}[runin]{\normalfont\sffamily\bfseries\color{Ink}}{}{0pt}{}[.]
\titlespacing*{\paragraph}{0pt}{0.9em}{0.6em}
\setcounter{secnumdepth}{3}
\setcounter{tocdepth}{2}

% ---------- Table des matières ----------
\titlecontents{section}[0pt]{\addvspace{0.7em}\sffamily\bfseries\color{Ink}}
  {}{}{\hfill\contentspage}[\addvspace{0.2em}]
\titlecontents{subsection}[2.2em]{\small\linespread{1}\selectfont}
  {\makebox[2.6em][l]{\color{Gris}\thecontentslabel}}{}{\titlerule*[0.7pc]{.}\contentspage}

% ---------- En-têtes et pieds de page ----------
\pagestyle{fancy}
\fancyhf{}
\renewcommand{\sectionmark}[1]{\markboth{#1}{}}
\fancyhead[L]{\footnotesize\sffamily\color{Gris}\leftmark}
\fancyhead[R]{\footnotesize\sffamily\scshape\color{Gris}Le Collé}
\fancyfoot[C]{\small\sffamily\color{Ink}\thepage}
\fancyfoot[L]{\scriptsize\sffamily\hyperlink{toc}{\color{Filet}sommaire}}
\renewcommand{\headrulewidth}{0.4pt}
\renewcommand{\headrule}{{\color{Filet}\hrule height\headrulewidth width\headwidth}}
\renewcommand{\footrulewidth}{0pt}
\fancypagestyle{plain}{\fancyhf{}\fancyfoot[C]{\small\sffamily\color{Ink}\thepage}\renewcommand{\headrulewidth}{0pt}}

% ---------- Énoncés ----------
\newtheoremstyle{enonce}{0pt}{0pt}{}{}{\sffamily\bfseries\color{Bleu}}{.}{0.5em}
  {\thmname{#1}\thmnumber{~#2}\thmnote{\ \normalfont\sffamily\color{Ink}(#3)}}
\newtheoremstyle{exemple}{0.9em}{0.9em}{}{}{\sffamily\bfseries\color{BleuMoyen}}{.}{0.5em}
  {\thmname{#1}\thmnumber{~#2}\thmnote{\ \normalfont\sffamily\color{Ink}(#3)}}
\newtheoremstyle{remarque}{0.7em}{0.7em}{}{}{\sffamily\itshape}{.}{0.5em}{\thmname{#1}\thmnote{\ (#3)}}
\theoremstyle{enonce}
\newtheorem{definition}{Définition}[section]
\newtheorem{theoreme}[definition]{Théorème}
\newtheorem{proposition}[definition]{Proposition}
\newtheorem{corollaire}[definition]{Corollaire}
\theoremstyle{exemple}
\newtheorem{exemple}[definition]{Exemple}
\theoremstyle{remarque}
\newtheorem*{remarque}{Remarque}

\tcbset{
  defi/.style={enhanced,breakable,colback=white,colframe=white,boxrule=0pt,arc=0pt,
     borderline west={1.2pt}{0pt}{Bleu},left=11pt,right=4pt,top=3pt,bottom=3pt,
     before skip=1.1em,after skip=1.1em},
  thm/.style={enhanced,breakable,colback=Fond,colframe=Fond,boxrule=0pt,arc=0pt,
     borderline west={2.4pt}{0pt}{Bleu},left=11pt,right=8pt,top=6pt,bottom=6pt,
     before skip=1.2em,after skip=1.2em}}
\tcolorboxenvironment{definition}{defi}
\tcolorboxenvironment{theoreme}{thm}
\tcolorboxenvironment{proposition}{thm}
\tcolorboxenvironment{corollaire}{thm}

\renewcommand{\proofname}{Démonstration}
\renewcommand{\qedsymbol}{\textcolor{BleuClair}{$\blacksquare$}}
\makeatletter
\renewenvironment{proof}[1][\proofname]{\par\pushQED{\qed}\normalfont
  \topsep0.5em\relax\trivlist\item[\hskip\labelsep\sffamily\itshape\color{Gris}#1.]\ignorespaces}
  {\popQED\endtrivlist\@endpefalse}
\makeatother

% ---------- Méthode : filets haut et bas, titre en bandeau ----------
\makeatletter\newcommand{\methlabelfix}{\protected@edef\@currentlabel{\thetcbcounter}}\makeatother
\newtcolorbox[auto counter,number within=section]{methode}[2][]{%
  enhanced,unbreakable,colback=white,colframe=white,boxrule=0pt,arc=0pt,
  borderline north={0.9pt}{0pt}{Bleu},borderline south={0.9pt}{0pt}{Bleu},
  left=4pt,right=4pt,top=8pt,bottom=8pt,before skip=1.4em,after skip=1.4em,
  before upper={\methlabelfix\phantomsection\addcontentsline{met}{metentry}{\protect\makebox[2.6em][l]{\thetcbcounter}#2}%
     {\sffamily\bfseries\small\color{white}\colorbox{Bleu}{\strut\ MÉTHODE~\thetcbcounter\ }}%
     \ \ {\sffamily\bfseries\color{Bleu}#2}\par\smallskip},#1}

% ---------- Repères pédagogiques : un filet de couleur, un mot-clé ----------
\newcommand{\definirbloc}[3]{% nom, couleur, libellé
  \newtcolorbox{#1}[1][]{enhanced,breakable,blanker,
    borderline west={1.4pt}{0pt}{#2},left=10pt,right=2pt,top=1pt,bottom=1pt,
    before skip=1em,after skip=1em,
    before upper={{\sffamily\bfseries\footnotesize\color{#2}\textls[120]{\MakeUppercase{#3}}}\quad},##1}}
\definirbloc{reflexe}{Bleu}{Réflexe}
\definirbloc{astuce}{Ocre}{Astuce}
\definirbloc{attention}{Brique}{Attention}
\definirbloc{enpratique}{BleuMoyen}{En pratique}
\definirbloc{plusloin}{Gris}{Pour aller plus loin}
\definirbloc{strategie}{Bleu}{Stratégie}
\definirbloc{lecture}{BleuMoyen}{Lire le schéma}
% compatibilité avec les parties reprises
\newenvironment{garde}{\begin{attention}}{\end{attention}}
\newenvironment{complementHP}{\begin{plusloin}}{\end{plusloin}}
\newenvironment{interpretation}{\begin{enpratique}}{\end{enpratique}}

% ---------- Exercices, corrigés ----------
\newcounter{exo}[section]
\renewcommand{\theexo}{\thesection.\arabic{exo}}
\newcommand{\niveau}[1]{{\color{Ocre}\footnotesize
  \ifcase#1\or$\bullet\circ\circ\circ\circ$\or$\bullet\bullet\circ\circ\circ$\or$\bullet\bullet\bullet\circ\circ$%
  \or$\bullet\bullet\bullet\bullet\circ$\or$\bullet\bullet\bullet\bullet\bullet$\fi}}
\newcommand{\serie}[1]{\par\needspace{8\baselineskip}\vspace{1.5em}%
  \noindent{\sffamily\bfseries\color{Bleu}#1}\par\nobreak\vspace{-0.4em}%
  \noindent{\color{Filet}\rule{\linewidth}{0.4pt}}\par\nobreak\vspace{0.1em}}
\newenvironment{exo}[3]{%
  \par\needspace{5\baselineskip}\vspace{0.75em}\refstepcounter{exo}\label{e:#1}%
  \def\exolab{#1}%
  \noindent{\sffamily\bfseries\color{Bleu}Exercice~\theexo}\quad\niveau{#2}%
  \hfill{\footnotesize\sffamily\color{Gris}#3}\par\nobreak\vspace{0.15em}}%
  {\par\nopagebreak\noindent\hfill{\scriptsize\sffamily\hyperref[c:\exolab]{\color{Gris}corrigé p.~\pageref{c:\exolab}}}\par}
\newenvironment{corrige}[1]{%
  \par\needspace{6\baselineskip}\vspace{1em}\phantomsection\label{c:#1}%
  \noindent{\sffamily\bfseries\color{Bleu}Exercice~\ref{e:#1}}\quad
  {\scriptsize\sffamily\hyperref[e:#1]{\color{Gris}énoncé p.~\pageref{e:#1}}}\par\nobreak\vspace{0.2em}}%
  {\hfill\textcolor{BleuClair}{$\square$}\par}
\newcommand{\etape}[1]{\par\smallskip\noindent{\sffamily\bfseries\small\color{BleuMoyen}#1.}\ }
\newenvironment{solution}{\par\noindent{\sffamily\bfseries\small\color{BleuMoyen}Solution.}\ }{\hfill\textcolor{BleuClair}{$\square$}\par}
\newcommand{\chercher}{\par\smallskip\noindent{\color{Filet}\rule[0.5ex]{1.2cm}{0.4pt}}\
  {\footnotesize\sffamily\itshape\color{Gris}laisser chercher l'étudiant avant de poursuivre}\par\smallskip}

% ---------- Guide de séance ----------

\newcommand{\marge}[1]{\leavevmode\marginpar{\raggedleft\scriptsize\sffamily\color{Gris}\textls[100]{\MakeUppercase{#1}}}\ignorespaces}
\newcommand{\Etape}[4]{%
  \subsection[Étape #1 --- #3 {\color{Gris}(#2)}]{#3}\label{etape:#1}%
  \noindent{\small\sffamily\color{Gris}{\color{Ocre}\bfseries Étape #1}\quad·\quad #2\quad·\quad #4}\par\medskip}
\newcommand{\question}[1]{\par\smallskip\noindent{\sffamily\bfseries\small\color{BleuMoyen}Question.}\ \textit{#1}\par\smallskip}
\newcommand{\rappel}[1]{\par\smallskip\noindent\phantomsection\addcontentsline{rap}{rapentry}{#1}%
  {\sffamily\bfseries\scriptsize\colorbox{Fond}{\color{Gris}\strut RAPPEL}}\ {\sffamily\small\color{Gris}#1.}\ }
\newcommand{\pourquoi}{\par\noindent{\sffamily\bfseries\small\color{Ocre}Pourquoi ?}\ }

% Ouverture de chapitre : \ouverture{objectifs}{prérequis}{durée}
\newcommand{\ouverture}[3]{%
  \begin{tcolorbox}[enhanced,colback=Fond,colframe=Fond,arc=0pt,boxrule=0pt,left=12pt,right=12pt,top=10pt,bottom=10pt,
     before skip=1.2em,after skip=1.6em]
  \begin{minipage}[t]{0.56\linewidth}
  {\sffamily\bfseries\small\color{Bleu}OBJECTIFS}\par\vspace{0.3em}\small #1
  \end{minipage}\hfill
  \begin{minipage}[t]{0.38\linewidth}
  {\sffamily\bfseries\small\color{Bleu}PRÉREQUIS}\par\vspace{0.3em}\small #2
  \ifx\relax#3\relax\else\par\vspace{0.6em}{\sffamily\bfseries\small\color{Bleu}DURÉE}\par\vspace{0.2em}\small #3\fi
  \end{minipage}
  \end{tcolorbox}}
% Fin de chapitre
\newenvironment{essentiel}{\par\needspace{12\baselineskip}\vspace{1.5em}%
  \begin{tcolorbox}[enhanced,unbreakable,colback=white,colframe=white,boxrule=0pt,arc=0pt,
   borderline north={1.4pt}{0pt}{Ocre},left=2pt,right=2pt,top=8pt,bottom=4pt]
  {\sffamily\bfseries\color{Ocre}L'ESSENTIEL DU CHAPITRE}\par\medskip}{\end{tcolorbox}}

% ---------- Notations ----------
\newcommand{\K}{\mathbb{K}}
\newcommand{\R}{\mathbb{R}}
\newcommand{\C}{\mathbb{C}}
\newcommand{\N}{\mathbb{N}}
\newcommand{\Z}{\mathbb{Z}}
\newcommand{\Mn}[1]{\mathcal{M}_{n}(#1)}
\newcommand{\GLn}[1]{\mathrm{GL}_n(#1)}
\newcommand{\LE}{\mathcal{L}(E)}
\newcommand{\rg}{\operatorname{rg}}
\newcommand{\vect}{\operatorname{Vect}}
\newcommand{\Mat}{\operatorname{Mat}}
\newcommand{\Id}{\operatorname{Id}}
\newcommand{\sh}{\operatorname{sh}}
\newcommand{\ch}{\operatorname{ch}}
\renewcommand{\th}{\operatorname{th}}
\DeclareMathOperator{\argsh}{argsh}
\newcommand{\vv}[1]{\overrightarrow{#1}}
\newcommand{\ve}[1]{\vec{#1}}
\newcommand{\ient}[2]{\llbracket #1,#2\rrbracket}
\newcommand{\dt}[1]{\left|\begin{matrix}#1\end{matrix}\right|}
\newcommand{\dis}{\displaystyle}

% ---------- Figures ----------
\pgfplotsset{trig format plots=rad}
\pgfplotsset{courbe/.style={axis lines=middle,axis equal image,
  axis line style={Gris,-{Stealth[length=5pt]}},
  xlabel={$x$},ylabel={$y$},ticklabel style={font=\scriptsize\sffamily,color=Gris},
  every axis x label/.style={at={(ticklabel* cs:1)},anchor=west,font=\small},
  every axis y label/.style={at={(ticklabel* cs:1)},anchor=south,font=\small},
  samples=300,smooth,width=0.6\linewidth}}
\tikzset{arcetude/.style={Bleu,line width=1.25pt},
         reste/.style={BleuClair,line width=1pt},
         tgte/.style={Ocre,line width=0.9pt},
         asy/.style={Gris,dashed,line width=0.6pt},
         pt/.style={circle,fill=Ink,inner sep=1.2pt}}

% ---------- Navigation ----------
\makeatletter
\newcommand*\l@rapentry{\@dottedtocline{1}{0em}{0em}}
\newcommand*\l@metentry{\@dottedtocline{1}{0em}{0em}}
\newcommand{\listemethodes}{\@starttoc{met}}
\newcommand{\listerappels}{\@starttoc{rap}}
\makeatother
% onglets de chapitre sur le bord droit
\usepackage{eso-pic}
\newif\ifonglets \ongletsfalse
\newcommand{\dessineronglet}{%
  \edef\ongcol{\ongletcouleur}%
  \edef\ongnum{\grandnum}%
  \pgfmathsetmacro{\ongy}{3.2+2.3*(\ongnum-1)}%
  \begin{tikzpicture}[remember picture,overlay]
  \fill[\ongcol] ([xshift=-0.75cm,yshift=-\ongy cm]current page.north east)
     rectangle ([yshift=-\ongy cm-1.9cm]current page.north east);
  \node[text=white,font=\sffamily\bfseries\large] at ([xshift=-0.37cm,yshift=-\ongy cm-0.95cm]current page.north east)
     {\ongnum};
  \end{tikzpicture}}
\AddToShipoutPictureBG{\ifonglets\ifnum\value{section}>0 \dessineronglet\fi\fi}
% blocs de navigation et d'enseignement
\definirbloc{commencer}{Gris}{Si vous commencez ici}
\definirbloc{horsprogramme}{Brique}{Hors programme}
\newcommand{\HP}{\,{\sffamily\bfseries\scriptsize\textcolor{Brique}{[HP]}}}
\newcommand{\source}[1]{{\footnotesize\sffamily\color{Gris}\,(#1)}}
\newcommand{\coupdepouce}[1]{\par\vspace{-0.2em}\noindent{\footnotesize\sffamily\color{Gris}%
  \textbf{Coup de pouce} (à donner si l'étudiant bloque) : #1}\par}
\newcommand{\pointarret}[1]{\par\vspace{1em}\noindent{\color{Filet}\rule[0.5ex]{\linewidth}{0.4pt}}\par\nobreak
  \noindent{\footnotesize\sffamily\color{Ocre}\bfseries POINT D'ARRÊT POSSIBLE}\quad{\footnotesize\sffamily\color{Gris}#1}\par\vspace{0.6em}}
\newcommand{\demarche}[1]{\noindent{\sffamily\bfseries\small\color{BleuMoyen}Démarche.}\ \textit{#1}\par\smallskip}

% ---------- Corrigés : idées, astuces, pièges, à retenir ----------
\newcommand{\choix}[1]{\par\noindent{\sffamily\bfseries\small\color{BleuMoyen}Choix de la méthode.}\ #1\par\smallskip}
\newcommand{\gain}[1]{\par\noindent{\sffamily\bfseries\small\color{Ocre}Astuce.}\ #1\par}
\newcommand{\piege}[1]{\par\noindent{\sffamily\bfseries\small\color{Brique}Piège.}\ #1\par}
\newtcolorbox{retenirbox}{enhanced,breakable,colback=Fond,colframe=Fond,arc=0pt,boxrule=0pt,
  borderline west={2pt}{0pt}{Ocre},left=8pt,right=6pt,top=4pt,bottom=4pt,before skip=0.7em,after skip=0.5em}
\newcommand{\retenir}[1]{\begin{retenirbox}{\sffamily\bfseries\footnotesize\color{Ocre}À RETENIR}\quad #1\end{retenirbox}}
\newcommand{\TD}[1]{{\color{Ocre}\bfseries TD\,#1}\ ·\ }
\newcommand{\q}[1]{\par\smallskip\noindent{\sffamily\bfseries\color{Bleu}#1}\ }
\providecommand{\iddots}{\mathinner{\mkern1mu\raise1pt\hbox{.}\mkern2mu\raise4pt\hbox{.}\mkern2mu\raise7pt\hbox{.}\mkern1mu}}


% =====================================================================
%  Entraînement — Réduction : Banque PT 2018, Mathématiques A, problème d'algèbre
%  Une source, deux PDF : le guide de séance (par défaut) et le corrigé (\def\CORRIGE{}).
%  Le sujet officiel n'est pas reproduit (droits d'auteur) : il est consultable sur WikiPrépa.
% =====================================================================
\providecommand{\Ker}{\operatorname{Ker}}\renewcommand{\Im}{\operatorname{Im}}
\providecommand{\Id}{\mathrm{Id}}
\fancyfoot[L]{}\fancyhead[L]{\footnotesize\sffamily\color{Gris}Entraînement · Réduction · Banque PT 2018 et 2015, Maths A · corrigé Le Collé}
\newcommand{\Partie}[1]{\par\needspace{5\baselineskip}\vspace{1.1em}\noindent
  {\sffamily\bfseries\large\color{Bleu}#1}\par\vspace{2pt}{\color{Filet}\rule{\linewidth}{0.4pt}}\par\smallskip}
\newcommand{\Q}[1]{\par\needspace{3\baselineskip}\medskip\noindent{\sffamily\bfseries\color{Bleu}#1}\enspace}
\newcommand{\Demande}[1]{{\small\color{Gris}\emph{#1}}\par\smallskip}
\newcommand{\lienDL}[1]{\par\noindent{\sffamily\bfseries\small\color{Ocre}Lien avec le DL.}\ #1\par}
\newcommand{\Ind}[2]{\par\noindent{\sffamily\bfseries\small\color{BleuMoyen}#1}\ #2\par}
\newcommand{\WP}{https://wikiprepa.fr/concours/ecrit/banque-pt_maths-a_pt_2018}
\newcommand{\WPsujet}{https://wikiprepa.fr/concours/data/banque-pt_maths-a_pt_2018.pdf}
\newcommand{\WPb}{https://wikiprepa.fr/concours/ecrit/banque-pt_maths-a_pt_2015}
\newcommand{\WPbsujet}{https://wikiprepa.fr/concours/data/banque-pt_maths-a_pt_2015.pdf}
\setlist[enumerate,1]{label=\arabic*.,leftmargin=1.6em,itemsep=3pt,topsep=3pt}

\begin{document}
\pagestyle{fancy}
\thispagestyle{plain}

\begin{tikzpicture}[remember picture,overlay]
\fill[Bleu] (current page.north west) rectangle ([yshift=-4.2cm]current page.north east);
\node[anchor=north west,text=white,font=\sffamily\footnotesize,inner sep=0] at ([xshift=2.2cm,yshift=-1.1cm]current page.north west)
  {\textls[120]{MATHÉMATIQUES · CPGE PT · SUJET DE CONCOURS PT}};
\node[anchor=north west,text=white,font=\sffamily\bfseries\LARGE,inner sep=0] at ([xshift=2.2cm,yshift=-1.75cm]current page.north west)
  {Réduction des endomorphismes et des matrices};
\node[anchor=north west,text=white,font=\sffamily\normalsize,inner sep=0] at ([xshift=2.2cm,yshift=-2.75cm]current page.north west)
  {Banque PT 2018 et 2015 · Mathématiques A · corrigé Le Collé (non officiel)};
\node[anchor=north east,text=white,font=\sffamily\scshape\large,inner sep=0] at ([xshift=-2.2cm,yshift=-1.05cm]current page.north east)
  {Le Collé};
\node[anchor=north east,text=white,font=\sffamily\small,inner sep=0] at ([xshift=-2.2cm,yshift=-2.8cm]current page.north east)
  {Séance du samedi 21 novembre 2026};
\end{tikzpicture}
\vspace*{2.9cm}

\begin{tcolorbox}[enhanced,colback=Fond,colframe=Fond,arc=0pt,boxrule=0pt,left=10pt,right=10pt,top=8pt,bottom=8pt]
\small
\textbf{\sffamily\color{Bleu}Les sujets.} Deux épreuves de \emph{Mathématiques A} de la Banque PT (filière PT, 4 heures,
calculatrice interdite), dont on traite le \textbf{problème d'algèbre linéaire} ; les exercices de probabilités sont écartés.
\begin{itemize}[leftmargin=1.2em,itemsep=1pt,topsep=2pt]
\item \textbf{Session 2018}, parties I à III : diagonalisation, $\R$ ou $\C$, polynômes annulateurs, sous-espaces stables
(\href{\WPsujet}{sujet en PDF}, \href{\WP}{page du sujet}). La partie II utilise les rotations de $\R^3$.
\item \textbf{Session 2015}, parties I à III : familles $(x,f(x),f^2(x),\dots)$, sous-espace engendré par les itérés, relation
polynomiale, sous-espaces stables d'un endomorphisme diagonalisable (\href{\WPbsujet}{sujet en PDF}, \href{\WPb}{page du sujet}).
\end{itemize}
\textbf{\sffamily\color{Bleu}Ce document.} Le sujet officiel n'est pas reproduit (droits d'auteur) : le lire d'abord sur WikiPrépa.
Ce corrigé est rédigé par Le Collé ; il n'est pas officiel. Des corrigés de professeurs (UPS) sont référencés sur les pages des
sujets. Chaque réponse commence par la méthode choisie et la raison de ce choix ; les calculs ont été vérifiés par calcul formel.
\end{tcolorbox}

% =====================================================================
%                              CORRIGÉ
% =====================================================================
On note $\Id$ pour $\Id_E$, et $E_\lambda(f)=\Ker(f-\lambda\Id)$.

\section*{Banque PT 2018 — Mathématiques A}\addcontentsline{toc}{section}{Banque PT 2018}\label{s:2018}

\Partie{Partie I — Deux matrices de $\mathcal M_3(\R)$}

\Q{I.1} \choix{pour savoir si une matrice est diagonalisable, on calcule $\chi$ et les dimensions des sous-espaces propres.
Trois valeurs propres distinctes en dimension 3 suffisent (corollaire du cours).}

$\chi_A(x)=\det(xI_3-A)=(x+1)(x-1)(x-2)$ : trois valeurs propres distinctes, \textbf{$A$ est diagonalisable dans $\R$}.

$\chi_B(x)=(x-4)(x-1)^2$ (développer selon la troisième ligne, qui vaut $(0,0,x-1)$). La valeur propre $1$ est double ;
$B-I_3=\begin{pmatrix}3&0&-3\\3&0&-3\\0&0&0\end{pmatrix}$ est de rang 1, donc $\dim E_1(B)=2=m(1)$. Avec $\dim E_4(B)=1$, la somme
des dimensions vaut 3 : \textbf{$B$ est diagonalisable dans $\R$}.
\piege{conclure « $B$ n'est pas diagonalisable car $\chi_B$ a une racine double ». Une racine double n'empêche rien : il faut
comparer $\dim E_1$ à la multiplicité.}

\Q{I.2} Le calcul donne \textbf{$A^2=B$}. (Contrôle : la première ligne de $A^2$ est $(1+2+1,\ 1+0-1,\ -1-1-1)=(4,0,-3)$.)

\Q{I.3} On résout $(A-\lambda I_3)X=0$ pour $\lambda=-1,1,2$ :
\[E_{-1}(A)=\vect\big((1,-1,1)\big),\quad E_1(A)=\vect\big((1,1,1)\big),\quad E_2(A)=\vect\big((1,1,0)\big).\]
\[P=\begin{pmatrix}1&1&1\\-1&1&1\\1&1&0\end{pmatrix},\qquad D=\begin{pmatrix}-1&0&0\\0&1&0\\0&0&2\end{pmatrix},\qquad A=PDP^{-1}.\]
\gain{vérifier chaque vecteur propre : $A(1,1,0)^T=(2,2,0)^T$, $A(1,1,1)^T=(1,1,1)^T$, $A(1,-1,1)^T=(-1,1,-1)^T$. C'est plus sûr
que de calculer $P^{-1}$.}

\Q{I.4} \choix{« sans calcul » : on utilise I.2 et I.3.}
$B=A^2=(PDP^{-1})(PDP^{-1})=PD^2P^{-1}$, et $D^2=\mathrm{diag}(1,1,4)$ est diagonale : \textbf{$B$ est diagonalisable}, avec la
même matrice de passage $P$. On retrouve les valeurs propres $1$ (double) et $4$.
\retenir{si $A=PDP^{-1}$, alors $Q(A)=P\,Q(D)\,P^{-1}$ pour tout polynôme $Q$ : un polynôme en une matrice diagonalisable est
diagonalisable, dans la même base.}

\Partie{Partie II — Rotations de $\R^3$}

\Q{II.1(a)} $f\circ f$ est la rotation d'axe $\R e_3$ et d'angle $\pi$ : c'est le \textbf{demi-tour d'axe $\R e_3$} (symétrie
orthogonale par rapport à cet axe) : $f^2(e_1)=-e_1$, $f^2(e_2)=-e_2$, $f^2(e_3)=e_3$.

\Q{II.1(b)} $f(e_1)=e_2$, $f(e_2)=-e_1$, $f(e_3)=e_3$ :
\[C=\begin{pmatrix}0&-1&0\\1&0&0\\0&0&1\end{pmatrix}.\]

\Q{II.1(c)} $\chi_C(x)=(x-1)(x^2+1)$. Dans $\R$, $\chi_C$ n'est pas scindé : \textbf{$C$ n'est pas diagonalisable dans $\R$}.
Dans $\C$, ses racines $1$, $i$, $-i$ sont distinctes : \textbf{$C$ est diagonalisable dans $\C$}. Enfin
$C^2=\mathrm{diag}(-1,-1,1)$ est déjà diagonale : \textbf{$C^2$ est diagonalisable dans $\R$ et dans $\C$}.
\retenir{une matrice réelle peut être diagonalisable dans $\C$ sans l'être dans $\R$ : il faut toujours préciser le corps.}

\Q{II.2(a)} $\|w\|=\sqrt{18}=3\sqrt2$, donc $\bar w=\frac1{3\sqrt2}(1,1,-4)$. On choisit $u$ unitaire orthogonal à $w$ :
$u=\frac1{\sqrt2}(1,-1,0)$. Pour que $(u,v,\bar w)$ soit orthonormée directe, on prend $v=\bar w\wedge u=-\frac13(2,2,1)$.
(Contrôle : $u\wedge v=\bar w$.) D'autres choix de $u$ sont possibles.

\Q{II.2(b)} Dans $\mathcal B'$, $g$ fixe $\bar w$ et tourne de $\frac\pi2$ dans le plan $(u,v)$ : $M_{\mathcal B'}=C$. Dans la base
canonique, avec $Q=(u\,|\,v\,|\,\bar w)$, orthogonale : $\Mat_{\mathcal B}(g)=QCQ^{-1}=QCQ^T$, soit
\[\Mat_{\mathcal B}(g)=\frac1{18}\begin{pmatrix}1&1&-4\\1&1&-4\\-4&-4&16\end{pmatrix}+\frac{1}{3\sqrt2}\begin{pmatrix}0&4&1\\-4&0&1\\-1&-1&0\end{pmatrix}.\]
(C'est la formule $\bar w\bar w^T+[\bar w]_\wedge$ d'une rotation d'angle $\frac\pi2$ ; le calcul de $QCQ^T$ la redonne.)

\Q{II.2(c)} $M_{\mathcal B'}=C$ et $M_{\mathcal B'}^2=C^2$ : d'après II.1(c), \textbf{$M_{\mathcal B'}$ n'est pas diagonalisable dans $\R$,
$M_{\mathcal B'}^2$ l'est}.
\retenir{$g^2$ diagonalisable n'entraîne pas $g$ diagonalisable. La partie III donne une hypothèse qui rend l'implication vraie :
les valeurs propres de $g^2$ strictement positives. Ici $g^2$ a la valeur propre $-1$.}

\Partie{Partie III — Polynômes annulateurs et diagonalisabilité}

\Q{III.1} Soit $y\in\Im g$ : $y=g(x)$, donc $f(y)=f\circ g(x)=0$. Ainsi $\Im g\subset\Ker f$.

\Q{III.2(a)} On développe : les deux produits valent $f^2-(\alpha+\beta)f+\alpha\beta\Id$. Les facteurs \emph{commutent}.

\Q{III.2(b)} Soit $v$ un vecteur propre, associé à $\lambda_j$. Grâce à (a), on peut mettre le facteur $(f-\lambda_j\Id)$ en
dernière position (le premier appliqué) : il envoie $v$ sur $0$, donc le produit aussi.

\Q{III.2(c)} \choix{« une base bien choisie » : $f$ est diagonalisable, donc il existe une base de vecteurs propres.}
Soit $(v_1,\dots,v_n)$ une base de vecteurs propres et $x=\sum x_kv_k$. L'application $u=(f-\lambda_1\Id)\circ\dots\circ(f-\lambda_p\Id)$
est linéaire et nulle sur chaque $v_k$ (b), donc $u(x)=\sum x_ku(v_k)=0$.
\retenir{si $f$ est diagonalisable de valeurs propres $\lambda_1,\dots,\lambda_p$, le polynôme $\prod(X-\lambda_k)$ est
annulateur de $f$. La suite montre une réciproque dans le cas de deux valeurs.}

\Q{III.3(a)} On cherche $a,b$ avec $(a+b)f-(a\alpha+b\beta)\Id=\Id$ : $a+b=0$ et $-(a\alpha+b\beta)=1$. D'où
\[a=\frac1{\beta-\alpha},\qquad b=-\frac1{\beta-\alpha}\qquad(\alpha\neq\beta).\]

\Q{III.3(b)} Pour tout $x$ : $x=a(f-\alpha\Id)(x)+b(f-\beta\Id)(x)=(f-\alpha\Id)(ax)+(f-\beta\Id)(bx)$. Donc
$E=\Im(f-\alpha\Id)+\Im(f-\beta\Id)$.

\Q{III.3(c)} $(\star)$ s'écrit $(f-\alpha\Id)\circ(f-\beta\Id)=0$ : par III.1, $\Im(f-\beta\Id)\subset\Ker(f-\alpha\Id)$. Les facteurs
commutent (III.2(a)), donc aussi $(f-\beta\Id)\circ(f-\alpha\Id)=0$ et $\Im(f-\alpha\Id)\subset\Ker(f-\beta\Id)$.

\Q{III.3(d)} Par (b) puis (c) : $E=\Im(f-\alpha\Id)+\Im(f-\beta\Id)\subset\Ker(f-\beta\Id)+\Ker(f-\alpha\Id)\subset E$.

\Q{III.3(e)} Si $x\in\Ker(f-\alpha\Id)\cap\Ker(f-\beta\Id)$, alors $\alpha x=f(x)=\beta x$, donc $(\alpha-\beta)x=0$ et $x=0$. La somme
est directe.

\Q{III.3(f)} En réunissant une base de $\Ker(f-\alpha\Id)$ et une base de $\Ker(f-\beta\Id)$ (l'un des deux peut être réduit à
$\{0\}$), on obtient une base de $E$ (base adaptée à la somme directe) formée de vecteurs propres : \textbf{$f$ est
diagonalisable}.
\piege{dire « $\alpha$ et $\beta$ sont les valeurs propres de $f$ ». On sait seulement que les valeurs propres sont parmi
$\alpha,\beta$ : si $f=\alpha\Id$, $\beta$ n'est pas valeur propre. Le raisonnement ci-dessus marche dans tous les cas.}

\Q{III.4(a)} \choix{un sous-espace propre d'un endomorphisme $h$ est stable par tout endomorphisme qui commute avec $h$.}
$f$ commute avec $f^2$. Si $x\in F_k$, alors $f^2(f(x))=f(f^2(x))=f(\lambda_kx)=\lambda_kf(x)$, donc $f(x)\in F_k$.

\Q{III.4(b)} Pour $x\in F_k$ : $f_k^2(x)=f^2(x)=\lambda_kx=\mu_k^2x$. Donc
$(f_k+\mu_k\Id)\circ(f_k-\mu_k\Id)=f_k^2-\mu_k^2\Id=0$ sur $F_k$.

\Q{III.4(c)} $-\mu_k\neq\mu_k$ car $\mu_k=\sqrt{\lambda_k}>0$ : c'est l'hypothèse de III.3, appliquée à $f_k$ dans l'espace
$F_k$ (ici l'hypothèse « valeurs propres strictement positives » sert). \textbf{$f_k$ est diagonalisable}, et III.3(e) donne
$F_k=F_k^+\oplus F_k^-$.

\Q{III.4(d)} $f^2$ est diagonalisable, donc $E=F_1\oplus\dots\oplus F_p$. Chaque $F_k=F_k^+\oplus F_k^-$, donc
$E=F_1^+\oplus F_1^-\oplus\dots\oplus F_p^+\oplus F_p^-$. Les vecteurs non nuls de $F_k^\pm$ sont des vecteurs propres de $f$ (de
valeur propre $\mp\mu_k$). En réunissant des bases de ces sous-espaces, on obtient une base de $E$ formée de vecteurs propres :
\textbf{$f$ est diagonalisable}.
\retenir{la démarche « diagonaliser par blocs » : sous-espaces stables, restriction, puis recollement des bases. C'est l'idée
la plus importante du sujet.}

\begin{essentiel}
\begin{itemize}[leftmargin=1.1em,itemsep=3pt]
\item Diagonalisable : comparer $\dim E_\lambda$ à la multiplicité ; $n$ valeurs propres distinctes suffisent.
\item $A=PDP^{-1}\Rightarrow A^2=PD^2P^{-1}$ ; plus généralement $Q(A)=PQ(D)P^{-1}$.
\item Diagonalisabilité dans $\R$ ou dans $\C$ : ce n'est pas la même question.
\item Si $(f-\alpha\Id)(f-\beta\Id)=0$ avec $\alpha\neq\beta$, alors $E=E_\alpha\oplus E_\beta$ et $f$ est diagonalisable.
\item Un sous-espace propre de $h$ est stable par tout $f$ qui commute avec $h$ ; on réduit alors la restriction.
\end{itemize}
\end{essentiel}


% =====================================================================
\clearpage
\section*{Banque PT 2015 — Mathématiques A}\label{s:2015}
\markboth{Banque PT 2015}{}
Ce problème prolonge directement le DL~3 : familles $(x,f(x),f^2(x),\dots)$, relation polynomiale, et le même argument
« un polynôme qui a trop de racines est nul ».

\Partie{Partie I — Un endomorphisme de $\R^4$}

\Q{I.1} \choix{une famille de trois vecteurs est liée dès qu'on trouve une relation non triviale : on cherche $f^2(e_1)$ comme
combinaison de $e_1$ et $f(e_1)$.}
$f(e_1)$ est la première colonne de $A$ : $f(e_1)=(-7,9,7,-4)$. Puis $f^2(e_1)=Af(e_1)=(-30,-90,-70,40)$. On constate que
$f^2(e_1)=-10f(e_1)-100e_1$ (première coordonnée : $70-100=-30$ ; les autres : $-90$, $-70$, $40$). Donc
$f^2(e_1)+10f(e_1)+100e_1=0$ : \textbf{la famille est liée}.
\gain{deviner la relation : $e_1$ n'a de composante que sur la première coordonnée, donc les trois dernières coordonnées de
$f^2(e_1)$ et de $f(e_1)$ doivent être proportionnelles : $(-90,-70,40)=-10\,(9,7,-4)$.}

\Q{I.2} $f(e_2)=(-16,-3,-4,-7)$ et $f^2(e_2)=(160,-70,40,70)$. De même, $f^2(e_2)=-10f(e_2)-100e_2$ : la famille est liée.

\Q{I.3} \choix{4 vecteurs en dimension 4 : un déterminant non nul suffit.}
\begin{align*}\det(e_1,f(e_1),e_2,f(e_2))&=\begin{vmatrix}1&-7&0&-16\\0&9&1&-3\\0&7&0&-4\\0&-4&0&-7\end{vmatrix}
=\begin{vmatrix}9&1&-3\\7&0&-4\\-4&0&-7\end{vmatrix}\\&=-1\times\begin{vmatrix}7&-4\\-4&-7\end{vmatrix}=-(-49-16)=65\neq0.\end{align*}
(Développement selon la première colonne, puis selon la deuxième colonne.) \textbf{$\mathcal B$ est une base de $\R^4$.}

\Q{I.4} \choix{une application linéaire nulle sur une base est nulle.}
Soit $g=f^2+10f+100\Id$, linéaire. $g(e_1)=0$ et $g(e_2)=0$ (questions 1 et 2). Comme $g$ commute avec $f$ (c'est un polynôme
en $f$), $g(f(e_1))=f(g(e_1))=0$ et de même $g(f(e_2))=0$. $g$ est nulle sur la base $\mathcal B$, donc $g=0$ : pour tout $x$,
$f^2(x)+10f(x)+100x=0$.

\Q{I.5} $f(e_1)$ est le 2\textsuperscript{e} vecteur de $\mathcal B$ ; $f(f(e_1))=-100e_1-10f(e_1)$ ; de même pour $e_2$ :
\[\Mat_{\mathcal B}(f)=\begin{pmatrix}0&-100&0&0\\1&-10&0&0\\0&0&0&-100\\0&0&1&-10\end{pmatrix}.\]
\lienDL{c'est une matrice « compagnon » par blocs, comme la matrice $R$ du DL dans une base cyclique.}

\Q{I.6} \choix{trouver les valeurs propres possibles grâce au polynôme annulateur.}
Si $f(x)=\lambda x$ avec $x\neq0$, la question 4 donne $(\lambda^2+10\lambda+100)x=0$, donc $\lambda^2+10\lambda+100=0$. Le
discriminant vaut $100-400<0$ : $f$ n'a \textbf{aucune valeur propre réelle}. \textbf{$A$ n'est pas diagonalisable dans $\R$.}
\retenir{toute valeur propre est racine de tout polynôme annulateur. Dans $\C$, les racines $-5\pm5i\sqrt3$ sont distinctes, et
le résultat de 2018 III.3 (deux racines distinctes) montre que $A$ est diagonalisable dans $\C$.}

\Partie{Partie II — Le sous-espace engendré par les itérés d'un vecteur}

\Q{II.1} Pour $y=\sum c_nx_n\in E_x$ (somme finie), $f(y)=\sum c_nx_{n+1}\in E_x$ par linéarité : $E_x$ est stable par $f$.

\Q{II.2} Par récurrence : $x_0=x\in F$, et si $x_n\in F$, alors $x_{n+1}=f(x_n)\in F$ car $F$ est stable. Tous les $x_n$ sont dans
le sous-espace $F$, donc aussi leurs combinaisons linéaires : $E_x\subset F$.
\retenir{$E_x$ est le plus petit sous-espace stable par $f$ qui contient $x$.}

\Q{II.3(a)} $(x_0)$ est libre car $x\neq0$, et une famille libre de $\R^d$ a au plus $d$ vecteurs. L'ensemble des entiers $k\geq1$
tels que $(x_0,\dots,x_{k-1})$ soit libre est non vide et majoré par $d$ : il a un plus grand élément $p$, avec $1\leq p\leq d$.

\Q{II.3(b)} Par maximalité de $p$, $(x_0,\dots,x_p)$ est liée, alors que $(x_0,\dots,x_{p-1})$ est libre : $x_p$ est combinaison
linéaire de $x_0,\dots,x_{p-1}$ (résultat du cours : on ajoute un vecteur à une famille libre et elle devient liée).
\piege{dire seulement « la famille est liée, donc $x_p$ est combinaison des autres ». Il faut la liberté de $(x_0,\dots,x_{p-1})$ :
dans une relation $\sum_{i\leq p}c_ix_i=0$ non triviale, $c_p\neq0$, sinon la relation porterait sur la famille libre.}

\Q{II.3(c)} Pour $i\leq p-2$, $f(x_i)=x_{i+1}\in E'_x$ ; et $f(x_{p-1})=x_p\in E'_x$ par (b). Les images des vecteurs générateurs
sont dans $E'_x$ : il est stable.

\Q{II.3(d)} $E'_x$ contient $x$ et est stable : par II.2, $E_x\subset E'_x$. L'inclusion inverse est claire. Donc $E_x=E'_x$, et
$\mathcal B_p$ engendre $E_x$ et est libre : c'est une base de $E_x$, et $\dim E_x=p$.

\Q{II.4} $\hat f(x_i)=x_{i+1}$ pour $i\leq p-2$ et $\hat f(x_{p-1})=\sum a_ix_i$ :
\[\Mat_{\mathcal B_p}(\hat f)=\begin{pmatrix}0&&&a_0\\1&\ddots&&a_1\\&\ddots&0&\vdots\\&&1&a_{p-1}\end{pmatrix}.\]

\Q{II.5} \choix{appliquer la relation au vecteur $x$ : chaque $\hat f^k(x)$ est un vecteur de la base.}
Si $\sum_{k=0}^{p-1}c_k\hat f^k=0$, on l'applique à $x_0$ : $\sum c_kx_k=0$, et $\mathcal B_p$ est libre, donc tous les $c_k$ sont nuls.

\Q{II.6(a)} On a $x_k=\hat f^k(x_0)$, et les puissances de $\hat f$ commutent :
$\hat f^p(x_k)=\hat f^k\big(\hat f^p(x_0)\big)=\hat f^k(x_p)=\hat f^k\Big(\sum_ia_ix_i\Big)=\sum_ia_i\hat f^k\big(\hat f^i(x_0)\big)=\sum_ia_i\hat f^i(x_k)$.

\Q{II.6(b)} L'endomorphisme $\hat f^p-\sum_ia_i\hat f^i$ est nul sur chaque $x_k$, donc sur la base $\mathcal B_p$ : il est nul.

\Partie{Partie III — Sous-espaces stables d'un endomorphisme diagonalisable}

\Q{III.1(a)} $f$ est diagonalisable, donc $E=E_1\oplus\dots\oplus E_p$ : tout $x$ s'écrit de façon \textbf{unique} $x=\sum x_i$
avec $x_i\in E_i$.

\Q{III.1(b)} Les $x_i$ ($i\leq q$) sont non nuls et sont des vecteurs propres associés à des valeurs propres distinctes : la
famille est libre (cours).

\Q{III.1(c)} $f(x_i)=\lambda_ix_i$, donc $f^k(x_i)=\lambda_i^kx_i$ et $f^k(x)=\sum_{i=1}^q\lambda_i^kx_i$ (pour tout $k$).

\Q{III.1(d)} \choix{c'est exactement le calcul du DL, III.3 : on regroupe selon les vecteurs propres.}
$\sum_{j=1}^q\alpha_jf^{j-1}(x)=\sum_{i=1}^q\Big(\sum_{j=1}^q\alpha_j\lambda_i^{j-1}\Big)x_i=\sum_{i=1}^qP(\lambda_i)x_i=0$. Par (b),
$P(\lambda_i)=0$ pour $i=1,\dots,q$.

\Q{III.1(e)} $P$ est de degré au plus $q-1$ et a $q$ racines distinctes : $P=0$, donc tous les $\alpha_j$ sont nuls. La famille
est libre.

\Q{III.1(f)} D'après (c), les $f^k(x)$ sont tous dans $V=\vect(x_1,\dots,x_q)$. $V$ contient $x$ et est stable par $f$ (engendré
par des vecteurs propres), donc $E_x\subset V$ (II.2). Par ailleurs $\vect(x,\dots,f^{q-1}(x))\subset E_x$, et ce sous-espace est de
dimension $q$ par (e), comme $V$. Donc $\vect(x,\dots,f^{q-1}(x))=E_x=V$. Comme $x_i=0$ pour $i>q$,
$V=\vect(x_1,\dots,x_p)$.

\Q{III.2} $F$ est stable et contient $x$ : $E_x\subset F$ (II.2). Or $x_i\in E_x$ (III.1(f)), donc $x_i\in F\cap E_i=F_i$.
\retenir{un sous-espace stable par un endomorphisme diagonalisable contient les « composantes propres » de chacun de ses
vecteurs : $F=\bigoplus_i(F\cap E_i)$.}

\Q{III.3(a)} $\chi_B(x)=x(x-1)(x-2)$ : trois valeurs propres distinctes en dimension 3, \textbf{$f$ est diagonalisable}.

\Q{III.3(b)} $E_0=\vect\big((1,-1,0)\big)$, $E_1=\vect\big((1,-1,-1)\big)$, $E_2=\vect\big((0,1,1)\big)$. (Contrôle :
$B(1,-1,-1)^T=(1,-1,-1)^T$ et $B(0,1,1)^T=(0,2,2)^T$.)

\Q{III.3(c)} \choix{par III.2, un sous-espace stable $F$ vérifie $F\subset F_0+F_1+F_2\subset F$, donc $F=F_0\oplus F_1\oplus F_2$ ;
chaque $F_i\subset E_i$ est une droite ou $\{0\}$.}
Chaque $F_i$ vaut $\{0\}$ ou $E_i$ ; réciproquement, toute somme de sous-espaces propres est stable. Les sous-espaces stables
sont donc les $2^3=8$ sous-espaces : $\{0\}$ ; les trois droites propres $E_0$, $E_1$, $E_2$ ; les trois plans $E_0\oplus E_1$,
$E_0\oplus E_2$, $E_1\oplus E_2$ ; et $\R^3$.
\piege{oublier des plans, ou en ajouter : un plan stable contient forcément deux droites propres.}

\begin{essentiel}
\begin{itemize}[leftmargin=1.1em,itemsep=3pt]
\item Une relation $P(f)=0$ donne les valeurs propres possibles : les racines de $P$.
\item $E_x=\vect(x,f(x),\dots)$ est le plus petit sous-espace stable contenant $x$ ; sa dimension est le premier $p$ où la famille
devient liée.
\item Regrouper selon les vecteurs propres, puis « un polynôme de degré $<q$ avec $q$ racines est nul » : le même argument que
dans le DL~3.
\item Sous-espaces stables d'un endomorphisme à valeurs propres simples : les sommes de droites propres.
\end{itemize}
\end{essentiel}

\end{document}
